% Quatre rayonnements placés sur l'échelle des domaines électromagnétiques (corrigé de l'exercice 1, chapitre 20)
\begin{tikzpicture}[x=8.5mm,y=1cm]
\def\hb{0.55}
\foreach \a/\b/\c/\n in {-11/-8/otNavy!22/{rayons X},-8/-6.4/violet!20/{UV},-6.1/-3/red!18/{infrarouge},-3/0/otOrange!22/{micro-ondes},0/3/otGreen!22/{ondes radio}}{%
  \fill[\c] (\a,0) rectangle (\b,\hb);
  \node[pcLegendeG] at ({(\a+\b)/2},{\hb/2}) {\n};}
\fill[otNavy!40] (-12,0) rectangle (-11,\hb);
\node[pcLegendeG] at (-11.5,{\hb/2}) {$\gamma$};
\shade[shading=otArcEnCiel] (-6.4,0) rectangle (-6.1,\hb);
\draw[otInk!60,line width=0.5pt] (-12,0) rectangle (3,\hb);
\draw[otInk!70,line width=0.5pt,-{Stealth[length=2mm]}] (-12,0) -- (3.5,0) node[pcLegende,anchor=north west] {$\lambda$ (m)};
\foreach \k in {-12,...,3}{\draw[otInk!70,line width=0.4pt] (\k,0) -- (\k,-0.08);}
\foreach \k/\t in {-12/{10^{-12}},-9/{10^{-9}},-6/{10^{-6}},-3/{10^{-3}},0/{1},3/{10^{3}}}{\node[pcLegende,anchor=north] at (\k,-0.08) {$\t$};}
\foreach \x/\t in {-10/{c : \qty{0,10}{nm}},-5/{d : \qty{10}{\micro m}},-0.91/{a et b : \qty{12}{cm}}}{%
  \draw[otRed!85!black,line width=0.8pt,-{Stealth[length=2mm]}] (\x,1.05) -- (\x,0.05);
  \node[pcLegendeG,anchor=south,text=otRed!85!black] at (\x,1.05) {\t};}
\end{tikzpicture}
